\documentclass[11pt]{article}
\usepackage[T1]{fontenc}
\usepackage{lmodern}
\usepackage{amsmath,amssymb,amsthm}
\usepackage{graphicx}
\usepackage[margin=0.88in]{geometry}
\usepackage[colorlinks=true,linkcolor=blue,citecolor=blue,urlcolor=blue]{hyperref}

\newtheorem{theorem}{Theorem}
\newtheorem{lemma}{Lemma}
\newtheorem{corollary}{Corollary}
\newtheorem{remark}{Remark}
\newcommand{\R}{\mathbb R}
\newcommand{\Gammaa}{\Gamma}
\newcommand{\dist}{\operatorname{dist}}

\title{Cone Covering Under Curvature Bounds\\at the Admissible Scale}
\author{Run \texttt{fa\_banach\_001}, agent \texttt{agent\_lane\_11}}
\date{17 August 2026}

\begin{document}
\maketitle

\section*{Source Question and Status}

Amenta and Kemppainen prove a cone-covering lemma for complete,
geometrically doubling Riemannian manifolds of nonnegative sectional
curvature.  They then write that they know no doubling Riemannian manifold
on which it fails and ask for a more precise determination of the spaces
admitting such coverings \cite[Remark~B.6]{AK}.

\begin{center}
\includegraphics[width=0.96\linewidth]{figures/open_question_crop.png}
\end{center}

This packet gives a broad sufficient condition allowing negative curvature:
sectional curvature may be as negative as a fixed multiple of
$-m(x)^{-2}$ on the ball of radius comparable to the admissibility scale
$m(x)$.  It includes the globally lower-bounded-curvature case whenever
$m$ is bounded, and therefore every compact Riemannian manifold with a
bounded admissibility function satisfying the source assumptions.  Status:
\emph{substantial partial result, likely valid}, subject to human review.

\section*{Setup and Proof Intuition}

Retain the source's complete geometrically doubling Riemannian manifold
$(X,g)$, weighted measure $\gamma$, and admissibility function
$m:X\to(0,\infty)$.  Thus $\gamma$ is doubling on every fixed enlargement
of the admissible balls and, for every $a>0$, there is $c_a$ such that
\begin{equation}\label{eq:C}
 d(x,y)\leq a m(x)\quad\Longrightarrow\quad m(x)\leq c_a m(y).
\end{equation}
Put $\beta=c_1$, as in the proof of source Lemma B.4, and fix the extension
$E^*$ supplied by source Lemma B.3 for this $\beta$.

Nonnegative curvature enters that proof only once: it supplies a finite net
of radial directions whose geodesics remain within $s/4$ of every other
radial geodesic.  Negative curvature makes radial geodesics separate, but
Rauch comparison quantifies the separation.  If curvature is controlled on
the scale $m(x)$, a uniformly finer angular net restores the same $s/4$
estimate.  The measure-theoretic part of the source proof is then unchanged.

\section*{Admissible-Scale Angular Comparison}

\begin{lemma}\label{lem:angular}
Suppose that for some $\kappa\geq0$ and a point $x\in X$,
\begin{equation}\label{eq:curvature-local}
 \sec_z\geq-\frac{\kappa^2}{m(x)^2}
 \qquad\text{for every }z\in \overline B(x,\beta m(x)).
\end{equation}
Set
\begin{equation}\label{eq:delta}
 \delta_{\kappa,\beta}=
 \begin{cases}
 \dfrac{\kappa\beta}{4\sinh(\kappa\beta)},&\kappa>0,\\[6pt]
 \dfrac14,&\kappa=0.
 \end{cases}
\end{equation}
If $v,w\in T_xX$ are unit vectors with
$\angle(v,w)\leq\delta_{\kappa,\beta}$, then their unit-speed geodesics
satisfy
\[
 d(\gamma_v(s),\gamma_w(s))\leq s/4
 \qquad(0\leq s\leq\beta m(x)).
\]
\end{lemma}

\begin{proof}
Join $v$ to $w$ by the short great-circle arc in the unit tangent sphere and
vary the corresponding radial geodesics.  The variation field is an
orthogonal Jacobi field $J$ with $J(0)=0$ and unit initial derivative when
the angular parameter is arclength.  Rauch comparison under
\eqref{eq:curvature-local} gives
\[
 |J(s)|\leq
 \begin{cases}
 \dfrac{m(x)}{\kappa}\sinh\!\left(\dfrac{\kappa s}{m(x)}\right),&\kappa>0,\\[6pt]
 s,&\kappa=0.
 \end{cases}
\]
The distance between the two endpoints is at most the length of this
angular variation.  For $q=\kappa s/m(x)\leq\kappa\beta$ and $\kappa>0$,
\[
 \frac{d(\gamma_v(s),\gamma_w(s))}{s}
 \leq \angle(v,w)\frac{\sinh q}{q}
 \leq \frac14\frac{\kappa\beta}{\sinh(\kappa\beta)}
                 \frac{\sinh q}{q}
 \leq\frac14,
\]
because $q\mapsto\sinh(q)/q$ is increasing.  The limiting case
$\kappa=0$ is immediate.
\end{proof}

\section*{Main Partial Result}

\begin{theorem}[cone covering at the curvature scale]\label{thm:main}
Assume the source hypotheses above and suppose that there is a fixed
$\kappa\geq0$ such that, for every $x\in X$,
\begin{equation}\label{eq:scaled-curvature}
 \sec_z\geq-\frac{\kappa^2}{m(x)^2}
 \qquad(z\in \overline B(x,\beta m(x))).
\end{equation}
Then the cone-covering conclusion of source Lemma B.4 holds: for every
bounded open $E\subset X$ and $x\in E$, there are
$x_1,\ldots,x_N\in X\setminus E$, with $N$ depending only on
$\dim X$, $\beta$, and $\kappa$, such that
\[
 \Gammaa(x)\setminus T(E^*)
 \subset\bigcup_{j=1}^N\Gammaa(x_j).
\]
Consequently, the source's pointwise estimate and its vector-valued
tent-space applications hold under \eqref{eq:scaled-curvature}.
\end{theorem}

\begin{proof}
Choose a $\delta_{\kappa,\beta}$-net
$v_1,\ldots,v_N$ in the unit sphere of $T_xX$.  Compactness of the
Euclidean unit sphere gives a cardinality bound depending only on
$\dim X$, $\beta$, and $\kappa$.  By Lemma~\ref{lem:angular}, every radial
unit geodesic $\gamma$ from $x$ has an index $j$ such that
\begin{equation}\label{eq:radial-control}
 d(\gamma(s),\gamma_{v_j}(s))\leq s/4
 \qquad(0\leq s\leq\beta m(x)).
\end{equation}

If $E=X$, then the left-hand side of the claimed covering is empty and the
empty family suffices.  Otherwise, for each $j$, let $t_j>0$ be the first
time at which
$\overline B(\gamma_{v_j}(t_j),t_j/4)$ meets $X\setminus E$, and choose
\[
 x_j\in(X\setminus E)\cap
 \overline B(\gamma_{v_j}(t_j),t_j/4).
\]
These times exist because $E$ is bounded.  Let
$(y,t)\in\Gammaa(x)\setminus T(E^*)$, write $r=d(x,y)<t$, and choose a
minimizing unit geodesic $\gamma$ from $x$ to $y$.  Pick $j$ satisfying
\eqref{eq:radial-control}.

First suppose $t_j\leq\beta m(x)$.  If $r>t_j$, then
\eqref{eq:radial-control} gives
\[
 d(y,x_j)\leq(r-t_j)+t_j/4+t_j/4<r<t.
\]
If $r\leq t_j$, then \eqref{eq:radial-control} shows that $y$ belongs to
the sector
\[
 R(v_j,t_j)=\bigcup_{0\leq s\leq t_j}
 B(\gamma_{v_j}(s),s/4),
\]
which is contained in $E$ by minimality of $t_j$.  Source Lemma B.3 gives
$B(y,2t_j)\subset E^*$.  Since
$(y,t)\notin T(E^*)$ means $B(y,t)\not\subset E^*$, one has $2t_j<t$.
The triangle estimate from the source proof gives $d(y,x_j)\leq2t_j<t$.

It remains to rule out $t_j>\beta m(x)$.  Since $r<t<m(y)$, applying
\eqref{eq:C} with the roles of $x,y$ reversed gives
$t<m(y)\leq c_1m(x)=\beta m(x)$.  Equation
\eqref{eq:radial-control} puts
$y\in R(v_j,\beta m(x))\subset E$.  Source Lemma B.3 therefore gives
$B(y,2\beta m(x))\subset E^*$, while
$B(y,t)\not\subset E^*$.  This forces $\beta m(x)<t$, a contradiction.
Thus one of the selected $x_j$ satisfies $d(y,x_j)<t$, proving the desired
cover.
\end{proof}

\begin{corollary}\label{cor:global}
Suppose $\sec X\geq-K$ globally and $m(x)\leq M$ for all $x$.  Then the
cone-covering property holds, with constants depending additionally on
$\sqrt K M$.  In particular, it holds on every compact Riemannian manifold
for every bounded admissibility function satisfying the source's local
doubling and comparability assumptions.
\end{corollary}

\begin{proof}
Take $\kappa=\sqrt K M$.  Since $m(x)\leq M$, one has
$-K\geq-\kappa^2/m(x)^2$, so \eqref{eq:scaled-curvature} follows from the
global lower curvature bound.  A compact Riemannian manifold has sectional
curvature bounded below.
\end{proof}

\section*{Scope, Upgrade Attempts, and Verification}

The theorem is a sufficient condition, not a characterization.  It rules
out compact manifolds and the usual bounded-admissibility, lower-curvature
settings as sources of counterexamples.  It also allows unbounded negative
curvature when $m(x)$ shrinks at least at the reciprocal curvature scale.
Rauch comparison genuinely uses sectional, rather than merely Ricci,
curvature.

An additional attempt to use geometric doubling alone did not close the
problem.  Doubling controls endpoint packings at each fixed scale, whereas
the proof requires a coherent radial net at all smaller scales; no such
implication or counterexample was found.  Hyperbolic space illustrates the
radial obstruction but is not doubling.

The run's cheap indexes contained no answer.  Searches through 17 August
2026 found the source and the later abstract use of the cone-covering
property in \cite{K}, but no counterexample or this admissible-scale
negative-curvature extension.  The companion script samples the exact
dimensionless inequality behind \eqref{eq:delta}; it is a consistency check,
not a substitute for Rauch comparison.

\section*{Human Review Recommendation}

Verify that the Rauch estimate is applicable to the full angular variation
inside $B(x,\beta m(x))$, including beyond the cut locus, and compare the
two case splits after \eqref{eq:radial-control} with source Lemma B.4.  These
are the only new geometric and bookkeeping steps.

\begin{thebibliography}{9}
\small
\bibitem{AK}
A. Amenta and M. Kemppainen,
\emph{Non-uniformly local tent spaces}, Publ. Mat. \textbf{59} (2015),
245--270; arXiv:1309.5758.

\bibitem{dC}
M. do Carmo, \emph{Riemannian Geometry}, Birkh\"auser, 1992; chapter on the
Rauch comparison theorem.

\bibitem{K}
M. Kemppainen, \emph{On vector-valued tent spaces and Hardy spaces
associated with non-negative self-adjoint operators}, Glasgow Math. J.
\textbf{58} (2016), 689--716; arXiv:1402.2886.
\end{thebibliography}

\end{document}
